\documentclass[11pt, a4paper]{article}
\input{preamble.tex}
\switchlang{ru}
\begin{document}

\title{1999 "--- IMicrosoft IntelliMouse Explorer mouse}
\date{}
\maketitle
\selectlanguage{russian}

Модель IntelliMouse Explorer 1.0, выпущенная в 1999 году, стала важной вехой в семействе Microsoft IntelliMouse "--- линейке мышей и трекболов, берущей начало с мыши 1995 года с колесом прокрутки. IntelliMouse Explorer была оснащена парой дополнительных кнопок (впервые в манипуляторах Microsoft) и кликабельным колесом, но её наиболее важной особенностью стала технология оптического отслеживания IntelliEye, которая произвела революцию в оптических мышах, избавив их от необходимости в специальном коврике. Встроив в это устройство ряд экспериментальных функций, Microsoft четко понимала их потенциал, рекламируя его как «самое радикальное достижение в области компьютерных технологий и дизайна» \cite{wiki}. На самом деле это утверждение нельзя считать совсем безосновательным, учитывая, что IntelliMouse Explorer стала родоначальником современных оптической мышей. Стоит отметить, что одновременно Microsoft выпустила устройство-близнец \emph{IntelliMouse with IntelliEye}, в котором практически идентичная аппаратная «начинка» была интегрирована в классический корпус образца 1995 года \cite{pcplayer, vogons}.

\begin{figure}[h]
    \centering
    \includegraphics[scale=0.6]{1999_microsoft_intellimouse_explorer/pic_60.jpg}
    \caption{Microsoft IntelliMouse Explorer mouse}
    \label{fig:MicrosoftIntelliMouseExplorerPic}
\end{figure}

\section*{Внешний вид}

Для IntelliMouse Explorer был разработан «футуристичный» по меркам 90-х годов серебристый корпус с серыми каплевидными вставками. Как отчетливо видно на рисунке \ref{fig:MicrosoftIntelliMouseExplorerTopBottom}, серебристое покрытие оказалось не слишком прочным и довольно быстро стиралось в местах сильного износа. По всей видимости, этот недостаток был следствием первой попытки, а не являлся запланированным поведением: по крайней мере, на более поздних моделях аналогичное покрытие выглядит лучше даже у экземпляров со значительным износом.

На верхней стороне расположены две основные кнопки (как у всех представителей семейства они крупные, занимают треть длины мыши), рассчитанные на нажатие указательным и средним пальцами, а между ними находится прорезиненное колесо прокрутки, при нажатии выполняющее функции средней кнопки мыши (рис. \ref{fig:MicrosoftIntelliMouseExplorerTopBottom}). Пара дополнительных кнопок расположена на левой стороне корпуса на одной из серых пластиковых вставок и рассчитана на нажатие большим пальцем.

\begin{figure}[h]
    \centering
    \includegraphics[scale=0.7]{1999_microsoft_intellimouse_explorer/top_30.jpg}
    \includegraphics[scale=0.7]{1999_microsoft_intellimouse_explorer/bottom_30.jpg}
    \caption{Microsoft IntelliMouse Explorer mouse, вид сверху и снизу}
    \label{fig:MicrosoftIntelliMouseExplorerTopBottom}
\end{figure}

Нижняя сторона мыши выполнена из прозрачного красного пластика, который заходит на передний к пользователю торец корпуса в виде еще одного каплевидного выступа. При движении мыши светодиод подсвечивает рабочую поверхность, что сопровождается красным свечением нижней части, наиболее заметной в упомянутом каплевидном выступе, играющем роль своеобразного хвостового огня. На нижней части корпуса расположены овальные накладки из низкофрикционного материала, облегчающие скольжение по рабочей поверхности, и полупрозрачный ярлык с техническими данными мыши.

\section*{Размеры и Эргономика}

Как можно видеть на рис. (рис. \ref{fig:MicrosoftPS2Size}), мышь достаточно крупная: ее корпус рассчитан на ладонь среднего или большого размера \cite{review}.

\begin{figure}[h]
    \centering
    \includegraphics[scale=0.6]{1999_microsoft_intellimouse_explorer/size_30.jpg}
    \caption{Изображение Microsoft IntelliMouse Explorer mouse на размерном коврике с шагом сетки 1~см}
    \label{fig:MicrosoftPS2Size}
\end{figure}

Учитывая высокую асимметрию корпуса мыши и отсутствие «левостороннего» варианта, устройства рассчитано исключительно на правшей. Однако к созданию удобной формы корпуса для них были приложены значительные усилия (сильный акцент на эргономике "--- тенденция, характерная для всего семейства IntelliMouse). Изогнутая форма корпуса позволяет комфортно опереться на неё ладонью \ref{fig:MicrosoftIntelliMouseExplorerHand} и нажимать кнопки при естественном положении кисти. В обзоре \cite{pcplayer} эргономика мыши характеризуется как превосходная, а в качестве незначительных недостатков упоминается, что устройство действительно удобно лежит только в довольно больших руках, и что до одной из двух боковых кнопок трудно дотянуться. Как можно наблюдать на рисунке \ref{fig:MicrosoftIntelliMouseExplorerHand}, показывающем положение модели женской руки на мыши, устройство нельзя назвать слишком большим, а проблема досягаемости, очевидно, касается кнопки, расположенной ближе всего к пользователю.

\begin{figure}[h]
    \centering
    \includegraphics[scale=0.6]{1999_microsoft_intellimouse_explorer/hand_30.jpg}
    \caption{Изображение Microsoft IntelliMouse Explorer mouse с моделью руки человека}
    \label{fig:MicrosoftIntelliMouseExplorerHand}
\end{figure}

Начиная с этой мыши появляется ограничение, касающееся невозможности работы на стекле и других сильно отражающих материалах, но при реальном использовании большинство обозревателей были полностью удовлетворены универсальностью технологии IntellyEye в отношении рабочих поверхностей \cite{review}. В то же время геймеры отмечали недостаточную точность этой мыши в жанрах, требующих быстрой реакции, и характеризовали её как отличное устройство для офисной работы и графики, но не рекомендуемое для динамичных игр \cite{pete}.

\section*{Подключение}

IntelliMouse Explorer стала первой мышью Microsoft с интерфейсом USB. Через переходник для этой модели также было доступно подключение к порту PS/2, а вот возможность работы через COM-порт уже отсутствовала (также впервые в линейке).

\section*{Внутреннее устройство}

Внутреннее устройство Microsoft IntelliMouse Explorer показано на рис. \ref{fig:MicrosoftPS2Inside}. С точки зрения аппаратного обеспечения, оптическая платформа IntellyEye очень похожа на ту, что используется в ручных сканерах, только с сильно уменьшенным полем зрения и очень высокой частотой обновления (в IntelliMouse Explorer 1.0 процессор с производительностью 16 MIPS обрабатывал 1500 крошечных кадров в секунду, а в последующих версиях мышей Microsoft для более плавной и отзывчивой работы частота кадров была увеличена в 4 раза). Разрешение CMOS-сенсора IntelliEye составляло около 400 DPI \cite{benchmark}.

\begin{figure}[h]
    \centering
    \includegraphics[width=\textwidth]{1999_microsoft_intellimouse_explorer/inside_30.jpg} 
    \caption{Microsoft IntelliMouse Explorer mouse в разобранном состоянии}
    \label{fig:MicrosoftPS2Inside}
\end{figure}

\begin{thebibliography}{9}
    \bibitem {wiki} IntelliMouse -- Wikipedia. \url{https://en.wikipedia.org/wiki/IntelliMouse}
    \bibitem {pcplayer} Microsoft IntelliMouse with IntelliEye // PC Player Computerspiele, 11/1999, p. 250. \url{https://archive.org/details/PC-Player-German-Magazine-1999-11/page/n137/mode/2up}
    \bibitem {vogons} Microsoft Intellimouse 1.0 with Intellieye USB (1999) \verb!\! VOGONS \url{https://www.vogons.org/viewtopic.php?t=78620}
    \bibitem {review} Tessler F. Microsoft IntelliMouse Explorer -- Macworld. Jan 31, 2000 \url{https://www.macworld.com/article/158316/intellimouse.html}
    \bibitem {pete} MS IntelliMouse With IntelliEye review. Eurogamer.net, 13 January 2000 \url{https://www.eurogamer.net/intellieye}
    \bibitem {benchmark} ESR -- ESReality MouseScore 2007. 24 December 2006 \url{https://www.esreality.com/post/1265679/esreality-mousescore-2007/}
\end{thebibliography}

\end{document}
